\documentclass[11pt,a4paper]{article}
\usepackage[top=2.3cm,bottom=2.3cm,left=2.2cm,right=2.2cm]{geometry}

\usepackage{amsmath,amssymb}
\usepackage{fontspec}
\setmainfont{Liberation Serif}
\setsansfont{Liberation Sans}
\setmonofont{DejaVu Sans Mono}[Scale=0.84]

\usepackage{xcolor}
\definecolor{brandblue}{RGB}{20, 55, 105}
\definecolor{accentblue}{RGB}{35, 95, 165}
\definecolor{codebg}{RGB}{248, 249, 251}
\definecolor{codeframe}{RGB}{215, 222, 230}
\definecolor{javakeyword}{RGB}{0, 45, 170}
\definecolor{javastring}{RGB}{5, 120, 20}
\definecolor{javacomment}{RGB}{120, 120, 120}
\definecolor{javanumber}{RGB}{20, 75, 220}

\definecolor{noteborder}{RGB}{30, 110, 65}
\definecolor{notebg}{RGB}{242, 249, 245}
\definecolor{warnborder}{RGB}{185, 75, 10}
\definecolor{warnbg}{RGB}{254, 248, 240}
\definecolor{tipborder}{RGB}{30, 95, 160}
\definecolor{tipbg}{RGB}{242, 247, 254}

\usepackage{fancyhdr}
\usepackage{lastpage}
\pagestyle{fancy}
\fancyhf{}
\lhead{\parbox[t]{0.58\textwidth}{\raggedright\small\sffamily\color{brandblue}\textbf{T03 --- Compilazione, Esecuzione e Deployment}}}
\rhead{\small\sffamily\color{gray}Programmazione II -- UniVR (2025/2026)}
\cfoot{\small\sffamily Pagina \thepage\ di \pageref{LastPage}}
\renewcommand{\headrulewidth}{0.5pt}
\renewcommand{\headrule}{\hbox to\headwidth{\color{brandblue!70!black}\leaders\hrule height \headrulewidth\hfill}}

\usepackage{titlesec}
\titleformat{\section}{\Large\bfseries\sffamily\color{brandblue}}{\thesection}{0.8em}{}[\titlerule]
\titleformat{\subsection}{\large\bfseries\sffamily\color{accentblue}}{\thesubsection}{0.8em}{}
\titleformat{\subsubsection}{\normalsize\bfseries\sffamily\color{brandblue!85!black}}{\thesubsubsection}{0.8em}{}

\usepackage{needspace}
\usepackage{fancyvrb}
\DefineVerbatimEnvironment{asciibox}{Verbatim}{
  fontsize=\fontsize{7.5pt}{9.2pt}\selectfont,
  frame=single,
  rulecolor=\color{codeframe},
  fillcolor=\color{codebg},
  framesep=2.5mm
}
\usepackage{listings}
\lstset{
  backgroundcolor=\color{codebg},
  frame=single,
  rulecolor=\color{codeframe},
  basicstyle=\ttfamily\footnotesize,
  keywordstyle=\color{javakeyword}\bfseries,
  stringstyle=\color{javastring},
  commentstyle=\color{javacomment}\itshape,
  numberstyle=\tiny\color{gray},
  numbers=left,
  numbersep=7pt,
  stepnumber=1,
  breaklines=true,
  breakatwhitespace=true,
  showstringspaces=false,
  tabsize=4,
  xleftmargin=12pt,
  framexleftmargin=12pt
}

\newcommand{\passthrough}[1]{#1}
\providecommand{\tightlist}{\setlength{\itemsep}{0pt}\setlength{\parskip}{0pt}}

\usepackage{tcolorbox}
\tcbuselibrary{skins,breakable}
\newtcolorbox{studynote}[1][Nota]{
  enhanced,
  breakable,
  colback=notebg,
  colframe=noteborder,
  fonttitle=\bfseries\sffamily\small,
  coltitle=white,
  title={#1},
  arc=2mm,
  boxrule=0.8pt,
  left=9pt, right=9pt, top=7pt, bottom=7pt
}

\newtcolorbox{studywarning}[1][Attenzione]{
  enhanced,
  breakable,
  colback=warnbg,
  colframe=warnborder,
  fonttitle=\bfseries\sffamily\small,
  coltitle=white,
  title={#1},
  arc=2mm,
  boxrule=0.8pt,
  left=9pt, right=9pt, top=7pt, bottom=7pt
}

\newtcolorbox{studytip}[1][Suggerimento]{
  enhanced,
  breakable,
  colback=tipbg,
  colframe=tipborder,
  fonttitle=\bfseries\sffamily\small,
  coltitle=white,
  title={#1},
  arc=2mm,
  boxrule=0.8pt,
  left=9pt, right=9pt, top=7pt, bottom=7pt
}

\usepackage{booktabs}
\usepackage{longtable}
\usepackage{array}
\usepackage{calc}

\usepackage{hyperref}
\hypersetup{
  colorlinks=true,
  linkcolor=accentblue,
  urlcolor=accentblue,
  citecolor=brandblue,
  pdfborder={0 0 0}
}

\title{\Huge\bfseries\sffamily\color{brandblue} T03 --- Compilazione, Esecuzione e Deployment \\[0.3em] \large\sffamily JVM, Bytecode, JDK vs JRE, Strumenti CLI (javac, java, jar) e Apache Maven}
\author{\textbf{Docente:} Prof. Michele Pasqua \quad---\quad \textbf{Appunti Revisione:} Wally}
\date{A.A. 2025/2026 -- Universit\`a degli Studi di Verona}

\begin{document}
\maketitle
\thispagestyle{fancy}
\vspace{-1.2em}
\noindent\textcolor{brandblue}{\rule{\textwidth}{0.8pt}}
\vspace{1.2em}

In questo blocco analizziamo l'ecosistema tecnologico di Java,
l'architettura della \textbf{Java Virtual Machine (JVM)}, i meccanismi
di compilazione, il \textbf{Dynamic Class Loading}, la gestione del
\textbf{Classpath} e le strategie di impacchettamento con file
\textbf{JAR} ed eseguibili
(\href{file:///home/wally/Documenti/GitHub/Programmazione-II/25-26/Teoria-e-Esercitazioni/Slides-20260902/T03\%20-\%20Building,\%20Running\%20and\%20Deployment\%20in\%20Java.pdf}{T03
- Building, Running and Deployment in Java.pdf}).

\begin{center}\rule{0.5\linewidth}{0.5pt}\end{center}

\subsubsection{1. Teoria: Concetti Teorici dalle
Slide}\label{teoria-concetti-teorici-dalle-slide}

\paragraph{A. Il Linguaggio Java e le sue Proprietà Fondamentali (Slide
1--11)}\label{a.-il-linguaggio-java-e-le-sue-proprietuxe0-fondamentali-slide-111}

\begin{itemize}
\tightlist
\item
  \textbf{Cenni Storici}:

  \begin{itemize}
  \tightlist
  \item
    Creato da \textbf{James Gosling} presso Sun Microsystems (progetto
    \emph{Oak} nel 1991, rilasciato ufficialmente come Java nel
    \textbf{1995}, poi acquisito da Oracle).
  \item
    Motto: \emph{``Write Once, Run Everywhere''} (WORA).
  \end{itemize}
\item
  \textbf{Le Caratteristiche Chiave del Linguaggio}:

  \begin{enumerate}
  \def\labelenumi{\arabic{enumi}.}
  \tightlist
  \item
    \textbf{Robustezza}:

    \begin{itemize}
    \tightlist
    \item
      \emph{Tipizzazione forte e statica (Strongly Typed)}: ogni
      variabile ed espressione ha un tipo noto a tempo di compilazione;
      i vincoli di tipo sono verificati dal compilatore prima
      dell'esecuzione.
    \item
      \emph{Nessuna manipolazione esplicita di puntatori}: in Java non
      esistono l'operatore di indirizzo \passthrough{\lstinline!\&!},
      l'aritmetica dei puntatori né la deallocazione manuale
      (\passthrough{\lstinline!free()!}). Questo elimina alla radice
      \emph{dangling pointers}, \emph{segmentation fault} e \emph{buffer
      overflow}.
    \item
      \emph{Controlli a runtime}: la JVM esegue verifiche automatiche su
      ogni accesso ad array
      (\passthrough{\lstinline!ArrayIndexOutOfBoundsException!}) e sui
      riferimenti (\passthrough{\lstinline!NullPointerException!}).
    \item
      \emph{Garbage Collection (GC) automatica}: un thread demone a
      bassa priorità della JVM individua ed elimina gli oggetti non più
      raggiungibili nello Heap, riducendo drasticamente i \emph{memory
      leak}.
    \item
      \emph{Gestione strutturata degli errori}: gestione degli stati
      anomali tramite eccezioni controllate e non controllate
      (\passthrough{\lstinline!try-catch-finally!}).
    \end{itemize}
  \item
    \textbf{Dinamicità}:

    \begin{itemize}
    \tightlist
    \item
      \emph{Caricamento e linking dinamico (Dynamic Linking)}: le classi
      non vengono collegate in un unico binario monolitico a
      compile-time, ma sono caricate in memoria dalla JVM
      \emph{on-demand}, solo quando effettivamente referenziate dal
      codice in esecuzione.
    \item
      \emph{Allocazione dinamica}: la dimensione delle strutture dati e
      degli array può essere determinata a runtime
      (\passthrough{\lstinline!new int[size]!}).
    \end{itemize}
  \item
    \textbf{Portabilità e Architettura a Bytecode}:

    \begin{itemize}
    \tightlist
    \item
      Nei linguaggi puramente compilati (C/C++), il compilatore genera
      direttamente codice macchina specifico per l'architettura target
      (x86, x64, ARM) e per il sistema operativo (Linux, Windows,
      macOS). Il programmatore deve ricompilare o mantenere build
      cross-platform separate.
    \item
      In Java, il compilatore (\passthrough{\lstinline!javac!}) traduce
      il sorgente \passthrough{\lstinline!.java!} in un formato
      intermedio indipendente dall'hardware: il \textbf{Bytecode}
      (\passthrough{\lstinline!.class!}).
    \item
      È la \textbf{JVM} a farsi carico della traduzione del bytecode in
      istruzioni macchina della CPU ospitante. La portabilità cessa di
      essere un onere del programmatore e diventa un servizio fornito
      dal runtime.
    \end{itemize}
  \item
    \textbf{Prestazioni e Compilatore JIT (\emph{Just-In-Time})}:

    \begin{itemize}
    \tightlist
    \item
      Java non è puramente interpretato. La JVM include un compilatore
      \textbf{JIT} che monitora a runtime le porzioni di bytecode
      eseguite più frequentemente (\emph{hot spots}) e le compila al
      volo in codice macchina nativo ottimizzato, memorizzandole nella
      cache del codice.
    \end{itemize}
  \end{enumerate}
\item
  \textbf{Java vs JavaScript}:

  \begin{itemize}
  \tightlist
  \item
    Citazione di Christian Heilmann: \emph{«Java is to JavaScript what
    Car is to Carpet»} (hanno in comune solo le prime quattro lettere).
    JavaScript nacque da Netscape con quel nome per motivi puramente
    commerciali; i due linguaggi non condividono né modello dei tipi, né
    macchina virtuale, né filosofia di design.
  \end{itemize}
\end{itemize}

\begin{center}\rule{0.5\linewidth}{0.5pt}\end{center}

\paragraph{B. Ecosistema Java: JDK vs JRE (Slide
11)}\label{b.-ecosistema-java-jdk-vs-jre-slide-11}

La distinzione tra ambiente di sviluppo e ambiente di esecuzione: -
\textbf{JDK (\emph{Java Development Kit})}: ambiente per programmatori:
- \passthrough{\lstinline!javac!}: il compilatore che traduce sorgenti
\passthrough{\lstinline!.java!} in bytecode
\passthrough{\lstinline!.class!}. - \passthrough{\lstinline!java!}: il
launcher applicativo che avvia la JVM. -
\passthrough{\lstinline!javadoc!}: generatore di documentazione HTML a
partire dai commenti speciali \passthrough{\lstinline!/** ... */!}. -
\passthrough{\lstinline!javap!}: disassemblatore di bytecode (permette
di ispezionare il bytecode generato). - \passthrough{\lstinline!jdb!}:
debugger da riga di comando. - \textbf{JRE (\emph{Java Runtime
Environment})}: ambiente minimale per l'utente finale: - \textbf{JVM
(\emph{Java Virtual Machine})}: il motore di esecuzione. - Librerie
standard delle API Java (il modulo base
\passthrough{\lstinline!java.base!} contenente
\passthrough{\lstinline!java.lang!},
\passthrough{\lstinline!java.util!}, \passthrough{\lstinline!java.io!},
ecc.). - Compilatore JIT integrato.

\begin{center}\rule{0.5\linewidth}{0.5pt}\end{center}

\paragraph{C. Compilazione, Esecuzione e Dynamic Class Loading (Slide
12--16)}\label{c.-compilazione-esecuzione-e-dynamic-class-loading-slide-1216}

\begin{enumerate}
\def\labelenumi{\arabic{enumi}.}
\item
  \textbf{La Struttura del Programma Minimo}:

\begin{lstlisting}[language=Java]
// Salvato obbligatoriamente in MyClass.java (stesso nome della classe pubblica)
public class MyClass {
    public static void main(String[] args) {
        System.out.println("Hello World!");
    }
}
\end{lstlisting}

  \begin{itemize}
  \tightlist
  \item
    \emph{Regola aurea}: In un file \passthrough{\lstinline!.java!} può
    essere presente al massimo \textbf{una sola classe pubblica}, e il
    nome del file deve coincidere esattamente (incluso il
    maiuscolo/minuscolo) con il nome di tale classe.
  \item
    \emph{Firma del metodo \passthrough{\lstinline!main!}}:

    \begin{itemize}
    \tightlist
    \item
      \passthrough{\lstinline!public!}: deve essere accessibile
      dall'esterno da parte del launcher della JVM.
    \item
      \passthrough{\lstinline!static!}: la JVM lo invoca direttamente
      sulla classe, senza dover prima istanziare un oggetto
      (\passthrough{\lstinline!MyClass obj = new MyClass()!}).
    \item
      \passthrough{\lstinline!void!}: non restituisce un codice numerico
      di ritorno (l'uscita anomala si gestisce con
      \passthrough{\lstinline!System.exit(code)!}).
    \item
      \passthrough{\lstinline!String[] args!}: array di stringhe
      contenente i parametri passati da riga di comando.
    \end{itemize}
  \end{itemize}
\item
  \textbf{Pipeline di Esecuzione della JVM}:
  \[\text{Sorgente } (.java) \xrightarrow{\texttt{javac}} \text{Bytecode } (.class) \xrightarrow{\text{Loader}} \text{Bytecode Verifier} \xrightarrow{\text{Interpreter / JIT}} \text{OS / Hardware}\]

  \begin{itemize}
  \tightlist
  \item
    \textbf{Bytecode Verifier}: componente di sicurezza fondamentale
    della JVM. Prima di eseguire qualsiasi file
    \passthrough{\lstinline!.class!}, scansiona le istruzioni per
    verificare che non violino i limiti dello stack, non convertano
    puntatori illegalmente e rispettino le regole di visibilità.
  \end{itemize}
\item
  \textbf{Dynamic Class Loading e Classpath
  (\passthrough{\lstinline!-cp!})}:

  \begin{itemize}
  \item
    La JVM non cerca i file nel filesystem a caso: utilizza il
    \textbf{Classpath} (una lista ordinata di directory e archivi
    \passthrough{\lstinline!.jar!}).
  \item
    Quando durante l'esecuzione il codice fa riferimento per la prima
    volta a una classe \passthrough{\lstinline!X!}:

    \begin{enumerate}
    \def\labelenumii{\arabic{enumii}.}
    \tightlist
    \item
      Il \passthrough{\lstinline!ClassLoader!} cerca
      \passthrough{\lstinline!X.class!} nella prima cartella specificata
      nel Classpath.
    \item
      Se non la trova, passa alla successiva.
    \item
      Se la trova, la carica in memoria e ne inizializza le strutture
      statiche.
    \item
      Se non la trova in nessuna cartella del Classpath, solleva
      \passthrough{\lstinline!ClassNotFoundException!} o
      \passthrough{\lstinline!NoClassDefFoundError!}.
    \end{enumerate}
  \item
    Esempio di esecuzione esplicita:

\begin{lstlisting}[language=bash]
java -cp . MyClass
\end{lstlisting}

    Il parametro \passthrough{\lstinline!-cp .!} istruisce la JVM a
    cercare le classi compilate nella directory corrente
    (\passthrough{\lstinline!.!}). Le classi standard di sistema
    (\passthrough{\lstinline!System!}, \passthrough{\lstinline!String!})
    vengono invece caricate automaticamente dal bootstrap classloader
    della JDK.
  \end{itemize}
\end{enumerate}

\begin{center}\rule{0.5\linewidth}{0.5pt}\end{center}

\paragraph{D. Deployment e Gestione dei File JAR (Slide
17--18)}\label{d.-deployment-e-gestione-dei-file-jar-slide-1718}

\begin{itemize}
\item
  \textbf{Anatomia di un file JAR (\emph{Java Archive})}:

  \begin{itemize}
  \tightlist
  \item
    Un file \passthrough{\lstinline!.jar!} è fisicamente un
    \textbf{archivio compresso in formato standard ZIP}.
  \item
    Raggruppa decine o centinaia di file
    \passthrough{\lstinline!.class!}, risorse (immagini, file di
    configurazione) e metadati.
  \end{itemize}
\item
  \textbf{Creazione manuale da CLI}:

\begin{lstlisting}[language=bash]
jar cvf MyJar.jar MyClass.class
\end{lstlisting}

  \begin{itemize}
  \tightlist
  \item
    \passthrough{\lstinline!c!}: \emph{create} (crea nuovo archivio).
  \item
    \passthrough{\lstinline!v!}: \emph{verbose} (mostra i file
    aggiunti).
  \item
    \passthrough{\lstinline!f!}: \emph{file} (il parametro successivo è
    il nome del file \passthrough{\lstinline!.jar!} da produrre).
  \end{itemize}
\item
  \textbf{Esecuzione tramite Classpath}:

\begin{lstlisting}[language=bash]
java -cp MyJar.jar MyClass
\end{lstlisting}
\item
  \textbf{JAR Auto-Eseguibile (\passthrough{\lstinline!java -jar!})}:

  \begin{itemize}
  \item
    Per consentire all'utente di lanciare l'archivio direttamente con
    \passthrough{\lstinline!java -jar MyJar.jar!}, il JAR deve
    specificare quale classe ospita il metodo
    \passthrough{\lstinline!main!}.
  \item
    Questa informazione risiede nel file
    \passthrough{\lstinline!META-INF/MANIFEST.MF!} sotto la direttiva:

\begin{lstlisting}
Main-Class: MyClass
\end{lstlisting}
  \item
    Creazione con file manifest esterno:

\begin{lstlisting}[language=bash]
jar cvfm MyClass.jar manifest.txt MyClass.class
\end{lstlisting}
  \item
    Creazione con sintassi moderna senza file di testo intermedio (flag
    \passthrough{\lstinline!e!} per entrypoint):

\begin{lstlisting}[language=bash]
jar cfe MyClass.jar MyClass MyClass.class
\end{lstlisting}
  \end{itemize}
\end{itemize}

\begin{center}\rule{0.5\linewidth}{0.5pt}\end{center}

\paragraph{E. Strumenti di Sviluppo: Da Editor a IntelliJ IDEA (Slide
19--26)}\label{e.-strumenti-di-sviluppo-da-editor-a-intellij-idea-slide-1926}

\begin{itemize}
\tightlist
\item
  \textbf{CLI vs IDE}: per progetti complessi, l'uso del terminale e di
  un editor di testo semplice diventa poco scalabile.
\item
  \textbf{Vantaggi dell'IDE (\emph{Integrated Development
  Environment})}:

  \begin{itemize}
  \tightlist
  \item
    \emph{Syntax highlighting} e \emph{live code inspection/linting}
    (segnalazione immediata degli errori di tipo e code smell prima
    ancora di compilare).
  \item
    \emph{Debugger integrato} con breakpoint, ispezione dello stack dei
    frame e monitoraggio variabili.
  \item
    \emph{Integrazione con Version Control (Git)} e \emph{Build
    Automation Tools (Maven / Gradle)}.
  \item
    \emph{Monito del docente}: l'auto-completamento o la generazione
    automatica di codice da parte dell'IDE è utile in ambito
    professionale, ma va limitata durante la fase di apprendimento per
    comprendere a fondo le regole sintattiche e semantiche del
    linguaggio.
  \end{itemize}
\item
  \textbf{Scelta del corso}: \textbf{IntelliJ IDEA} (gratuito in
  versione Community o con licenza Education per studenti universitari).
\end{itemize}

\begin{center}\rule{0.5\linewidth}{0.5pt}\end{center}

\subsubsection{\texorpdfstring{2. Codice: Collegamento Pratico con il
Tuo Progetto: Maven e
\texttt{guida-jar.md}}{2. Codice: Collegamento Pratico con il Tuo Progetto: Maven e guida-jar.md}}\label{codice-collegamento-pratico-con-il-tuo-progetto-maven-e-guida-jar.md}

Nel tuo workspace
\href{file:///home/wally/Documenti/GitHub/Programmazione-II/25-26/esercizi-wally-25-26}{\passthrough{\lstinline!esercizi-wally-25-26!}},
hai formalizzato questi esatti concetti in modo eccellente in
\href{file:///home/wally/Documenti/GitHub/Programmazione-II/25-26/esercizi-wally-25-26/guida-jar.md}{\passthrough{\lstinline!guida-jar.md!}}
e nei file di configurazione Maven:

\begin{enumerate}
\def\labelenumi{\arabic{enumi}.}
\tightlist
\item
  \textbf{Il Fully Qualified Name e il Filesystem}:

  \begin{itemize}
  \tightlist
  \item
    Come visto a lezione, se una classe appartiene al package
    \passthrough{\lstinline!es01!}, il bytecode generato deve risiedere
    fisicamente nel percorso \passthrough{\lstinline!es01/Date.class!}.
  \item
    Quando si invoca la JVM, il nome della classe da passare deve essere
    il FQN (\passthrough{\lstinline!es01.MainDate!}).
  \end{itemize}
\item
  \textbf{Automazione del Packaging con Maven}:

  \begin{itemize}
  \tightlist
  \item
    Invece di lanciare manualmente \passthrough{\lstinline!jar cfe!},
    nei singoli moduli (es.
    \href{file:///home/wally/Documenti/GitHub/Programmazione-II/25-26/esercizi-wally-25-26/es01/pom.xml}{\passthrough{\lstinline!es01/pom.xml!}})
    usi il plugin
    \textbf{\passthrough{\lstinline!maven-assembly-plugin!}}:
  \end{itemize}

\begin{lstlisting}[language=XML]
<plugin>
    <groupId>org.apache.maven.plugins</groupId>
    <artifactId>maven-assembly-plugin</artifactId>
    <version>3.7.1</version>
    <configuration>
        <archive>
            <manifest>
                <mainClass>es01.MainDate</mainClass>
            </manifest>
        </archive>
        <descriptorRefs>
            <descriptorRef>jar-with-dependencies</descriptorRef>
        </descriptorRefs>
    </configuration>
    <executions>
        <execution>
            <id>make-assembly</id>
            <phase>package</phase>
            <goals>
                <goal>single</goal>
            </goals>
        </execution>
    </executions>
</plugin>
\end{lstlisting}

  Questo genera automaticamente nella cartella
  \passthrough{\lstinline!target/!} il cosiddetto \emph{Fat JAR}
  contenente sia il bytecode dell'applicazione sia il file
  \passthrough{\lstinline!META-INF/MANIFEST.MF!} precompilato.
\end{enumerate}

\begin{center}\rule{0.5\linewidth}{0.5pt}\end{center}

\subsubsection{3. Test: Verifica dei Comandi da
Terminale}\label{test-verifica-dei-comandi-da-terminale}

Puoi verificare l'intero ciclo di build ed esecuzione da riga di comando
testando il modulo \passthrough{\lstinline!es01!} tramite lo script:

\begin{lstlisting}[language=bash]
./esercizi.sh package es01
\end{lstlisting}

Lo script esegue internamente: 1.
\passthrough{\lstinline!mvn clean package -pl es01!} 2. Individua il
file JAR autoprodotto in
\passthrough{\lstinline!es01/target/es01-*-jar-with-dependencies.jar!}
3. Esegue \passthrough{\lstinline!java -jar ...!} verificando che
l'entrypoint \passthrough{\lstinline!es01.MainDate!} risponda
correttamente.

\begin{center}\rule{0.5\linewidth}{0.5pt}\end{center}

\begin{studynote}[Nota di Studio] Con T03 abbiamo concluso l'intera panoramica propedeutica
(Introduzione, Paradigmi/UML e Tooling/JVM).

A partire dal prossimo blocco entriamo nel codice sorgente vero e
proprio con la \textbf{Lezione 4 / Slide T04: \emph{Types and Objects}}
e il file ufficiale del docente
\href{file:///home/wally/Documenti/GitHub/Programmazione-II/25-26/Teoria-e-Esercitazioni/Codice-20260902/Lezione04/Example.java}{\passthrough{\lstinline!Example.java!}}:
- Tipi primitivi (dimensioni, range, valori di default) vs Reference
Types. - Rappresentazione dei caratteri (\passthrough{\lstinline!char!}
come intero senza segno a 16 bit, codifica ASCII e Unicode
\passthrough{\lstinline!\\u0056!}). - Regole di inizializzazione:
variabili di istanza/statiche (default zero/null) vs variabili locali
dello stack (nessun default \(\rightarrow\) errore a compile-time!). -
Operatori logici standard (\passthrough{\lstinline!\&!},
\passthrough{\lstinline!|!}) vs operatori a corto circuito
(\emph{short-circuit} \passthrough{\lstinline!\&\&!},
\passthrough{\lstinline!||!}) e relativi effetti collaterali
(\passthrough{\lstinline!k++!}). - Scope dei blocchi di codice annidati
(\passthrough{\lstinline!\{ ... \}!}) e shadowing/visibilità delle
variabili.
\end{studynote}

\textbf{Dammi conferma per aprire la Lezione 4 / T04 e analizzare teoria
e codice di \passthrough{\lstinline!Example.java!}!}


\end{document}
